%--------------------------------------------------------------------------------
%
%
%--------------------------------------------------------------------------------
\pagebreak
\section*{RSM — Reset System Mask}
\addcontentsline{toc}{section}{RSM — Reset System Mask}

\begin{iPageDescEntry}{Syntax}
    \texttt{RSM RegR}
\end{iPageDescEntry}

\begin{iPageDescEntry}{Format}
    The \texttt{RSM} instruction uses the following instruction format: \\

    \begin{flushleft}
    \drawInstrFormat
        {0,1,3,5,6.5,13,13.5,14,14.5,15,15.5,16}
        {\OpCodeGrpSYS, \OpCodeRSM, RegR, 0, 0, 0, 0, 0, 0, 0, 0}
    \end{flushleft}
\end{iPageDescEntry}

\begin{iPageDescEntry}{Description}
    The \texttt{RSM} (Reset System Mask) instruction clears bits in the system
    mask register that are encoded in the instruction. The previous mask values are returned in \texttt{RegR}.
\end{iPageDescEntry}

\begin{iPageDescEntry}{Operation}
    \texttt{SystemMask <- SystemMask \& instr.[5:0];}
\end{iPageDescEntry}

\begin{iPageDescEntry}{Exceptions}
    - Privileged instruction trap.
\end{iPageDescEntry}

\begin{iPageDescEntry}{Notes}
     None.
\end{iPageDescEntry}

// ?? what is RegR used for ? return old mask ?


